%-----------------------------------------------------------------------
%  Packages, palette, frame furniture and the file-pair frame template.
%  \input by main.tex -- not compiled on its own.
%  Shares no file with any other document under notes/, so this deck can be
%  its own Overleaf project and go to one reader alone.
%-----------------------------------------------------------------------
\usepackage[T1]{fontenc}
\usepackage[utf8]{inputenc}
\usepackage{lmodern}
\usepackage{amsmath}
\usepackage{booktabs}
\usepackage{array}
\usepackage{tikz}
\usetikzlibrary{arrows.meta,positioning,fit,backgrounds}
\usepackage{microtype}

%-------------------------------------------------- palette
% Four roles, used consistently on every frame: what upstream has today, what
% is missing, what this project changed, and why it is safe. Reading the deck
% by colour alone should still tell the story.
\definecolor{ImpBlue} {HTML}{1F4E79}   % structure, and "the change"
\definecolor{ImpSlate}{HTML}{5A6672}   % upstream, as it stands
\definecolor{ImpAmber}{HTML}{A85C00}   % the gap
\definecolor{ImpGreen}{HTML}{2E6B3E}   % why it is safe / verified
\definecolor{ImpRule} {HTML}{C8D0D8}
\definecolor{ImpWash} {HTML}{F2F5F8}

%-------------------------------------------------- theme furniture
\usetheme{default}
\usecolortheme{default}
\setbeamertemplate{navigation symbols}{}
\setbeamercolor{structure}{fg=ImpBlue}
\setbeamercolor{normal text}{fg=black!88}
\setbeamerfont{frametitle}{series=\bfseries,size=\large}
\setbeamerfont{framesubtitle}{size=\footnotesize}
\setbeamercolor{framesubtitle}{fg=ImpSlate}

% Frame title: colour, then a hairline. No shaded bar -- these frames carry
% coloured boxes of their own and a heavy title competes with them.
% \makeatletter: this file is \input after \documentclass, so @ is not a
% letter here and the \@empty test below would otherwise not parse.
\makeatletter
\setbeamertemplate{frametitle}{%
  \vspace{0.2em}%
  {\usebeamerfont{frametitle}\color{ImpBlue}\insertframetitle\par}%
  \ifx\insertframesubtitle\@empty\else
    {\usebeamerfont{framesubtitle}\color{ImpSlate}\insertframesubtitle\par}%
  \fi
  \vspace{2pt}{\color{ImpRule}\rule{\textwidth}{0.4pt}}\par\vspace{-0.4em}}
\makeatother

\setbeamertemplate{footline}{%
  \hbox{\begin{beamercolorbox}[wd=\paperwidth,ht=2.2ex,dp=4.0ex,leftskip=1.2em,rightskip=1.2em]{}%
    \scriptsize\color{ImpSlate}%
    Tapenade-native AD hooks for MITgcm\hfill\insertframenumber%
  \end{beamercolorbox}}}

\setbeamertemplate{itemize item}{\color{ImpBlue}\rule[0.28ex]{0.5em}{0.5pt}}
\setbeamertemplate{itemize subitem}{\color{ImpSlate}\rule[0.28ex]{0.35em}{0.4pt}}
\setlength{\leftmargini}{1.2em}
\setlength{\leftmarginii}{1.1em}


%-------------------------------------------------- text macros
% \code renders identifiers verbatim-ish without escaping every underscore by
% hand. \detokenize hands TeX the characters rather than the catcodes, which is
% the whole reason a deck full of Fortran names is writable at all. Do not put
% a backslash or a percent sign inside it.
\newcommand{\code}[1]{\texttt{\detokenize{#1}}}
% CPP directives need their own macro: beamer grabs a frame body as a
% macro argument, so a literal # in it is read as a parameter number and
% the frame dies at its \end. \char35 sidesteps the catcode entirely.
\newcommand{\cpp}[1]{\texttt{\char35\relax\detokenize{#1}}}
\newcommand{\file}[1]{\texttt{\small\detokenize{#1}}}

%-------------------------------------------------- the file-pair template
% Every file-pair frame opens with the same header: the upstream file, the file
% in this repository that shadows it, and the size of the diff between them.
% The diff size is the argument -- these are small, guarded additions, not
% rewrites -- so it is on the frame rather than in the notes.
% ---------------------------------------------- the file-pair block
% Each modified file is shown as: our file (linked to it on GitHub) -> the
% MITgcm file it shadows, then three aligned rows saying what checkpoint69m
% does, what was missing for Tapenade, and what we changed. The fourth
% question -- why it still fits MITgcm's structure -- is the same for every
% file in a group, so it is answered once per frame by \Safe.

% The URL survives the underscores because #1 is substituted into \href's
% first argument, which hyperref normalises. Verified: the emitted PDF
% annotations carry the exact paths.
\newcommand{\ghfile}[1]{%
  \href{https://github.com/tsgonerog/Proj_ImPACTS/blob/main/MITgcm_c69m/mysetups/DINO_1deg/#1}%
       {\color{ImpBlue}\underline{\texttt{\detokenize{#1}}}}}

% One labelled row inside a file block: fixed-width label, hanging indent so
% wrapped text lines up under the text rather than under the label.
\newcommand{\pfrow}[3]{%
  \par\vspace{0.25pt}\noindent
  \hangindent=5.0em\hangafter=1
  \makebox[5.0em][l]{\labelface{#1}{#2}}#3\par}

% #1 our file  #2 the MITgcm file it shadows  #3 size of the difference
% #4 what checkpoint69m does  #5 what was missing  #6 what we changed
\newcommand{\pairfile}[6]{%
  \par\vspace{3.5pt}
  \begin{minipage}[t]{0.80\textwidth}\scriptsize\raggedright
    \ghfile{#1}\ $\rightarrow$\ {\color{ImpSlate}\texttt{\detokenize{#2}}}
  \end{minipage}\hfill
  \begin{minipage}[t]{0.19\textwidth}\scriptsize\raggedleft
    {\color{ImpSlate}#3}
  \end{minipage}%
  \par\vspace{1pt}
  {\scriptsize
   \pfrow{ImpSlate}{c69m}{#4}%
   \pfrow{ImpAmber}{Missing}{#5}%
   \pfrow{ImpBlue}{Change}{#6}}}

% The four labelled sections of a file-pair frame. A run-in coloured label
% rather than a boxed block: four blocks on one frame is a wall, four labels
% reads as a list and leaves room for the words that matter.
% ONE label formula for the whole deck, so it cannot drift again.
%
% Latin Modern Sans ships no bold small-caps shape: T1/lmss/bx/sc does not
% exist, so \bfseries\scshape silently rendered as T1/lmss/bx/n -- bold
% upright, the small caps dropped -- and LaTeX raised "Font shape
% `T1/lmss/bx/sc' undefined" plus its end-of-run "Some font shapes were not
% available" summary. Bold + \MakeUppercase asks for a shape the family
% actually has, gives the label-like appearance the small caps were reaching
% for, and is what \lbl already used for every section label -- so the strips
% and the section labels now share one style instead of two.
\newcommand{\labelface}[2]{{\color{#1}\bfseries\MakeUppercase{#2}}}

\newcommand{\lbl}[2]{\par\vspace{2.5pt}{\scriptsize\labelface{#1}{#2}}\par\vspace{0.5pt}}
\newcommand{\Upstream}{\lbl{ImpSlate}{In MITgcm checkpoint69m}}
\newcommand{\Gap}     {\lbl{ImpAmber}{What is missing for Tapenade}}
\newcommand{\Change}  {\lbl{ImpBlue} {The change}}
\newcommand{\Why}     {\lbl{ImpBlue} {Why it is necessary}}

% The closing strip every frame ends on, in the same place every time. \Safe
% is the common case -- the non-interference claim in one sentence -- and
% \Strip is for the few frames whose closing point is a caveat or a scope
% note rather than a safety claim, so the label never says the wrong thing.
\newcommand{\Strip}[3]{%
  \vfill
  \begin{beamercolorbox}[wd=\textwidth,sep=4pt]{safebox}%
    \scriptsize\labelface{#1}{#2}\ \ #3
  \end{beamercolorbox}}
\newcommand{\Safe}[1]{\Strip{ImpGreen}{Upstream-safe}{#1}}
\setbeamercolor{safebox}{bg=ImpGreen!8,fg=black}

\newcommand{\vfin}{\code{viscFacInAd}}

\setlength{\parskip}{1.5pt}
